\subsection{Hardware validation setup}
\label{sec:method:hardware}

The simulation study deliberately omits aerodynamic drag, motor asymmetry, battery sag
and actuator wear, so its absolute errors are optimistic and only its \emph{relative}
comparisons are claimed. To test whether the ranking survives contact with a real
vehicle, the controller recommended by the benchmark is validated on hardware in a
programmable wind tunnel with motion-capture ground truth. The setup is described here;
the measurements are reported in Sec.~\ref{sec:results:hardware}.

\paragraph{Vehicle.}
A Crazyflie~2.1 Brushless --- the airframe the simulation is parameterized on
(\SI{43.4}{\gram}, thrust-to-weight $1.88$) --- carrying a Lighthouse positioning deck,
so that the onboard state estimate available in flight is the physical counterpart of
the measurement model of Sec.~\ref{sec:method:sensor}.
% TODO: state where each controller executes (onboard firmware vs. offboard at 100 Hz
% over the radio) and the payload attachment method, once the bench is assembled.

\paragraph{Programmable wind field.}
Wind is generated by a WindShaper fan-array wind
generator\cite{windshaper_datasheet,windshaper_dantec} installed at the Institute of
Mechanics, Materials and Civil Engineering, UCLouvain\cite{uclouvain_windshaper}. Unlike
a conventional closed-circuit tunnel, which produces one nominally uniform stream, the
WindShaper is an array of independently driven \emph{wind pixels}: each module of
approximately \SI{25}{\centi\meter} $\times$ \SI{25}{\centi\meter} carries a $3\times3$ grid of counter-rotating fan
pairs, and modules tile to the required test-section area. Each pixel is commanded
separately, so the generated field is a programmable function of space and time,
$u = f(x,y,t)$, with time steps of \SI{0.1}{\second} through a documented Python
interface; the reported envelope is a mean speed from \SI{2}{\meter\per\second} up to
approximately \SI{20}{\meter\per\second} in the open-jet configuration, ramp rates of
about \SI{4}{\meter\per\second\squared}, and turbulence intensity programmable over
roughly $5$--$30\,\%$.
% TODO: replace the envelope figures above with the numbers measured on the UCLouvain
% installation (module count, test-section dimensions, calibrated speed and turbulence
% range at the flight volume) before submission.

This capability is the reason the facility was chosen, and it is what makes the
hardware campaign a validation rather than a loosely analogous experiment. The two wind
scenarios used in simulation are defined by their \emph{structure}, not merely their
magnitude: a spatially uniform steady field, and a steady field with a periodic
component at \SI{0.7}{\hertz} plus broadband turbulence. Both are directly expressible
as WindShaper set points --- the first as a constant pixel command, the second as a
sinusoidal modulation superimposed on a turbulence-intensity setting --- so the physical
and simulated conditions correspond one-to-one, and the sim-to-real gap we report is a
property of the vehicle and controller rather than of a mismatched disturbance.

\paragraph{Ground truth.}
Vehicle pose is measured by an OptiTrack PrimeX~22 motion-capture
system\cite{optitrack_primex22}: $2048\times1088$ pixels at \SI{360}{\hertz} (up to
\SI{1000}{\hertz} at reduced vertical resolution), a $79^\circ\times49^\circ$ field of
view with the standard \SI{6.8}{\milli\meter} lens, and a typical three-dimensional
accuracy of \SI{\pm0.2}{\milli\meter} with hardware synchronization across cameras. At
roughly two orders of magnitude below the tracking errors under study, and at more than
three times the \SI{100}{\hertz} control rate, the system is an unambiguous reference
for both roles it plays here. First, it supplies the true position used to evaluate
exactly the same metric as in simulation --- $\mathrm{RMSE}_{\mathrm{3D}}$ over the same
Lissajous references with the same \SI{1}{\second} warm-up exclusion --- so simulated
and flown numbers are directly comparable. Second, it provides the reference against
which the onboard Lighthouse estimate is differenced in flight, which turns the sensor
model of Sec.~\ref{sec:method:sensor} into a testable object: the update rate, the
quasi-static bias and the noise levels we simulate are all motion-capture-referenced
quantities, and the flight data either reproduces them or does not.

\paragraph{Conditions and protocol.}
The flight matrix mirrors the simulated one within the constraints of the test volume:
the same Lissajous reference at the speeds the physical envelope admits, a steady wind
and a gust profile matched to the simulated force magnitudes, and a payload step of the
same \SI{23}{\percent} of vehicle weight. Each cell is flown repeatedly so that hardware
error bars are comparable in kind to the seed statistics of
Sec.~\ref{sec:method:stats}, and the warm-up exclusion, metric definition and reference
parameterization are unchanged from simulation.
% TODO: fix the repetition count and the achievable speed tiers once the volume and the
% tunnel calibration are known.

\paragraph{What the campaign is designed to falsify.}
Three simulation claims are stated in advance as the campaign's targets, so that a
negative outcome is informative: \emph{(i)}~that under degraded onboard sensing with
wind the three leading stacks are statistically indistinguishable
(Sec.~\ref{sec:results:deployment}); \emph{(ii)}~that the hybrid adaptive stack absorbs
a payload step so completely that its tracking error is unchanged from its own nominal
value; and \emph{(iii)}~that the variance behaviour predicted under the simulated
Lighthouse model --- occasional near-failures on unlucky bias draws rather than a
uniform degradation --- appears with a real positioning deck.
